\documentclass[a4paper,12pt]{article}

\usepackage[T2A]{fontenc}
\usepackage[utf8]{inputenc}
\usepackage[russian]{babel}
\usepackage{amsmath,amssymb,amsthm,mathtools}
\usepackage{helvet}
\usepackage{paratype}
\renewcommand{\familydefault}{\sfdefault}
\usepackage{icomma}
\usepackage[a4paper,left=19mm,right=19mm,top=18mm,bottom=20mm]{geometry}
\usepackage[nopatch=footnote]{microtype}
\usepackage{tikz}
\usetikzlibrary{calc,arrows.meta,positioning,shapes.geometric}
\usepackage{pgfplots}
\pgfplotsset{compat=1.18}
\usepackage{tabularx,booktabs}
\usepackage{enumitem}
\usepackage{hyperref}
\usepackage{parskip}
\usepackage{fancyhdr}
\usepackage{setspace}
\usepackage{xcolor}

\definecolor{accent}{HTML}{008080}
\definecolor{darkaccent}{HTML}{004D4D}
\definecolor{textgray}{HTML}{333333}
\definecolor{softgray}{HTML}{E9EEEE}

\hypersetup{
    colorlinks=true,
    linkcolor=darkaccent,
    urlcolor=darkaccent
}

\setlength{\parindent}{0pt}
\onehalfspacing

\setlist[itemize]{
    leftmargin=1.7em,
    itemsep=0.25em,
    topsep=0.35em
}

\setlist[enumerate]{
    leftmargin=1.9em,
    itemsep=0.35em,
    topsep=0.35em
}

\color{textgray}

\pagestyle{fancy}
\fancyhf{}
\fancyhead[L]{\small Марина Селиверстова}
\fancyhead[R]{\small ОГЭ: графики и параметры}
\fancyfoot[C]{\thepage}
\setlength{\headheight}{15pt}
\renewcommand{\headrulewidth}{0.3pt}
\renewcommand{\footrulewidth}{0pt}

\newcommand{\sectionrule}{%
    \par
    \vspace{-0.4em}
    {\color{accent}\rule{\linewidth}{0.5pt}}
    \vspace{0.4em}
}

\newcommand{\key}[1]{%
    \textcolor{darkaccent}{\textbf{#1}}
}

\tikzset{
    flowstart/.style={
        ellipse,
        draw=darkaccent,
        very thick,
        align=center,
        minimum width=36mm,
        minimum height=13mm,
        fill=white
    },
    flowproc/.style={
        rectangle,
        rounded corners=3pt,
        draw=darkaccent,
        very thick,
        align=center,
        text width=48mm,
        minimum height=15mm,
        fill=white
    },
    flowwide/.style={
        rectangle,
        rounded corners=3pt,
        draw=darkaccent,
        very thick,
        align=center,
        text width=58mm,
        minimum height=17mm,
        fill=white
    },
    flowdec/.style={
        diamond,
        aspect=2.15,
        draw=darkaccent,
        very thick,
        align=center,
        text width=34mm,
        inner sep=1.5pt,
        fill=white
    },
    flowline/.style={
        -{Stealth[length=3mm,width=2mm]},
        very thick,
        draw=accent
    },
    flowlabel/.style={
        font=\small,
        fill=white,
        inner sep=1.5pt,
        text=darkaccent
    }
}

\newcommand{\startflowchartlandscape}{%
    \clearpage
    \pdfpagewidth=420mm
    \pdfpageheight=297mm
    \newgeometry{left=18mm,right=18mm,top=16mm,bottom=16mm}%
    \setlength{\textwidth}{384mm}%
    \setlength{\textheight}{265mm}%
    \setlength{\hsize}{\textwidth}%
    \setlength{\linewidth}{\textwidth}%
    \setlength{\columnwidth}{\textwidth}%
}

\newcommand{\finishflowchartpage}{%
    \clearpage
    \pdfpagewidth=210mm
    \pdfpageheight=297mm
    \restoregeometry%
    \setlength{\linewidth}{\textwidth}%
}

\begin{document}

% ============================================================
% ТИТУЛЬНЫЙ ЛИСТ
% ============================================================

\begin{titlepage}
\thispagestyle{empty}
\centering
\vspace*{22mm}

{\color{darkaccent}\Large\bfseries Чек-лист\par}

\vspace{8mm}

{\Huge\bfseries
Рациональные функции, параболы,\par}

\vspace{2mm}

{\Huge\bfseries
выколотые точки и параметры\par}

\vspace{8mm}

{\Large ОГЭ по математике\par}

\vfill

{\large Марина Селиверстова\par}

\vspace{5mm}

{\large Дата: \today\par}

\vspace{12mm}

{\large Лёвин Артём Александрович\par}
{\large эксклюзивно для Марины Селиверстовой\par}

\vfill

\begin{tikzpicture}[remember picture,overlay]
\draw[
    accent!55,
    line width=0.9pt
]
    ([xshift=18mm,yshift=18mm]current page.south west)
    .. controls
    ([xshift=60mm,yshift=31mm]current page.south west)
    and
    ([xshift=-70mm,yshift=8mm]current page.south east)
    ..
    ([xshift=-18mm,yshift=20mm]current page.south east);
\end{tikzpicture}

\end{titlepage}

% ============================================================
% ОСНОВНОЙ ЧЕК-ЛИСТ
% ============================================================

\section*{1. Главная идея занятия}
\sectionrule

В задачах ОГЭ график сложной дробно-рациональной функции часто после сокращения превращается в знакомую параболу. Ключевой момент: \key{ограничения исходного выражения сохраняются}. Поэтому после построения упрощённого графика нужно удалить точки, соответствующие запрещённым значениям $x$.

Типовая конструкция:
\[
y=\frac{P(x)(x-a)}{a-x}.
\]

Поскольку
\[
a-x=-(x-a),
\]
при $x\ne a$ получаем
\[
y=-P(x),
\qquad
x\ne a.
\]

На графике функции $y=-P(x)$ точка с абсциссой $x=a$ должна быть \key{выколота}.

\section*{2. ОДЗ и сокращение алгебраической дроби}
\sectionrule

\key{Шаг 1.} Найдите ОДЗ \emph{до сокращения}. Знаменатель не может быть равен нулю.

\key{Шаг 2.} Разложите числитель и знаменатель на множители, если это возможно.

\key{Шаг 3.} Сокращайте только одинаковые множители. Помните:
\[
x-a=-(a-x).
\]

Например,
\[
\frac{x-2}{2-x}
=
\frac{x-2}{-(x-2)}
=
-1,
\qquad
x\ne2.
\]

\key{Шаг 4.} После сокращения обязательно перенесите исходное ограничение в запись упрощённой функции.

\medskip

\textbf{Типичная ошибка.}
После сокращения множителя $(x-a)$ значение $x=a$ не возвращается в область определения. Сокращение меняет форму записи, но не исходную ОДЗ.

\section*{3. Как построить параболу после упрощения}
\sectionrule

Пусть после преобразований получена функция
\[
y=ax^2+bx+c,
\qquad
a\ne0.
\]

Для уверенного построения достаточно пройти следующие пункты.

\begin{enumerate}

\item
\key{Вершина параболы:}
\[
x_0=-\frac{b}{2a},
\qquad
y_0=f(x_0).
\]

\item
\key{Пересечения с осью $Ox$:}
положите $y=0$ и решите квадратное уравнение
\[
ax^2+bx+c=0.
\]

\item
\key{Пересечение с осью $Oy$:}
положите $x=0$. Получите точку
\[
(0;c).
\]

\item
\key{Направление ветвей:}
\[
a>0
\Rightarrow
\text{ветви вверх},
\qquad
a<0
\Rightarrow
\text{ветви вниз}.
\]

\item
Если точек недостаточно, составьте небольшую таблицу. Выбирайте значения $x$, симметричные относительно $x_0$.

\item
В самом конце нанесите \key{выколотые точки}, заданные ограничениями исходной функции.

\end{enumerate}

Если $b=0$, то
\[
x_0=0,
\]
и осью симметрии параболы является ось $Oy$.

Тогда удобно брать пары
\[
x=\pm1,
\qquad
x=\pm2
\]
и так далее.

\section*{4. Общие точки графиков и параметр}
\sectionrule

\subsection*{4.1. Прямая $y=m$}

Прямая
\[
y=m
\]
горизонтальна.

Число общих точек определяется её положением относительно графика.

Для параболы важно учитывать:

\begin{itemize}
\item уровень вершины;
\item направление ветвей;
\item выколотые точки: пересечение, попавшее в выколотую точку, \emph{не считается}.
\end{itemize}

\subsection*{4.2. Прямая $y=kx$}

Любая прямая вида
\[
y=kx
\]
проходит через начало координат.

Чтобы найти общие точки с графиком $y=f(x)$, приравниваем формулы:
\[
f(x)=kx.
\]

Если после этого получается квадратное уравнение, его дискриминант управляет количеством алгебраических точек пересечения:
\[
\begin{array}{ccl}
D>0 &\Rightarrow& 2\text{ корня},\\
D=0 &\Rightarrow& 1\text{ корень},\\
D<0 &\Rightarrow& 0\text{ корней}.
\end{array}
\]

При наличии выколотой точки возникает \key{дополнительный случай}. Два корня квадратного уравнения могут соответствовать двум пересечениям упрощённых графиков, но одно из них может быть запрещено ОДЗ. Тогда у исходных графиков останется ровно одна общая точка.

Если выколота точка
\[
A(x_h;y_h),
\qquad
x_h\ne0,
\]
то прямая $y=kx$, проходящая через неё, имеет коэффициент
\[
k=\frac{y_h}{x_h}.
\]

Это значение нужно рассматривать отдельно от случая касания
\[
D=0.
\]

% ============================================================
% БЛОК-СХЕМА
% ============================================================

\startflowchartlandscape

\thispagestyle{empty}

\noindent
\begin{minipage}{\textwidth}
\centering

{\Large\bfseries\color{darkaccent}
Алгоритм: график после сокращения и поиск параметра
\par}

\vspace{7mm}

\begin{tikzpicture}[
    node distance=15mm and 18mm,
    font=\small
]

\node[flowstart] (start)
{Исходная\\дробно-рациональная функция};

\node[flowproc,right=of start] (odz)
{Найти ОДЗ\\по исходному знаменателю};

\node[flowproc,right=of odz] (simp)
{Разложить на множители\\и сократить};

\node[flowproc,right=of simp] (keep)
{Записать упрощённую функцию\\вместе с ограничениями};

\node[flowproc,right=of keep] (parab)
{Построить параболу:\\вершина, оси, дополнительные точки};

\node[flowproc,below=of parab] (holes)
{Удалить точки, запрещённые ОДЗ};

\node[flowdec,left=23mm of holes] (type)
{Какая прямая\\задана?};

\node[
    flowwide,
    below left=18mm and 25mm of type
] (mline)
{$y=m$: двигать горизонтальную прямую и считать только существующие пересечения};

\node[
    flowwide,
    below right=18mm and 25mm of type
] (kline)
{$y=kx$: приравнять формулы $f(x)=kx$};

\node[
    flowproc,
    below=of kline
] (disc)
{Для касания потребовать $D=0$\\и найти соответствующие $k$};

\node[
    flowproc,
    left=22mm of disc
] (holek)
{Если выколота $A(x_h;y_h)$ и $x_h\ne0$, проверить $k=y_h/x_h$};

\node[
    flowwide,
    below=19mm of $(disc)!0.5!(holek)$
] (verify)
{Объединить кандидаты и проверить ОДЗ:\\сколько общих точек осталось у исходных графиков?};

\node[
    flowstart,
    below=of verify
] (finish)
{Записать ответ};

\draw[flowline] (start) -- (odz);
\draw[flowline] (odz) -- (simp);
\draw[flowline] (simp) -- (keep);
\draw[flowline] (keep) -- (parab);
\draw[flowline] (parab) -- (holes);
\draw[flowline] (holes) -- (type);

\draw[flowline]
(type)
--
node[flowlabel,above left] {$y=m$}
(mline);

\draw[flowline]
(type)
--
node[flowlabel,above right] {$y=kx$}
(kline);

\draw[flowline]
(mline.south)
|-
(verify.west);

\draw[flowline]
(kline)
--
(disc);

\draw[flowline]
(kline.south west)
|-
(holek.north);

\draw[flowline]
(disc)
--
(verify);

\draw[flowline]
(holek)
--
(verify);

\draw[flowline]
(verify)
--
(finish);

\end{tikzpicture}

\end{minipage}

\finishflowchartpage

% ============================================================
% ПРИМЕР 1
% ============================================================

\section*{5. Разобранный пример: парабола с выколотой вершиной}
\sectionrule

Рассмотрим функцию
\[
y=
\frac{(x+1)(x^2-4x+3)}{x-1}.
\]

Сначала фиксируем ОДЗ:
\[
x\ne1.
\]

Разложим квадратный трёхчлен:
\[
x^2-4x+3
=
(x-1)(x-3).
\]

Тогда при $x\ne1$
\[
y=
\frac{(x+1)(x-1)(x-3)}{x-1}
=
(x+1)(x-3)
=
x^2-2x-3.
\]

Вершина:
\[
x_0=
-\frac{-2}{2}
=
1,
\]
\[
y_0=
1-2-3
=
-4.
\]

Корни:
\[
x^2-2x-3=0,
\]
\[
(x+1)(x-3)=0.
\]

Следовательно, точки пересечения с осью $Ox$:
\[
(-1;0),
\qquad
(3;0).
\]

С осью $Oy$:
\[
(0;-3).
\]

Значение $x=1$ запрещено, поэтому вершина
\[
(1;-4)
\]
\key{выколота}.

Для горизонтальной прямой $y=m$ получаем:

\begin{center}
\begin{tabular}{cc}
\toprule
Значение $m$ & Число общих точек \\
\midrule
$m<-4$ & $0$ \\
$m=-4$ & $0$ \\
$m>-4$ & $2$ \\
\bottomrule
\end{tabular}
\end{center}

% ============================================================
% ПРИМЕР 2
% ============================================================

\section*{6. Разобранный пример: когда $y=kx$ имеет ровно одну общую точку}
\sectionrule

Пусть
\[
y=
\frac{(x^2+1)(x-2)}{2-x}.
\]

ОДЗ:
\[
x\ne2.
\]

Так как
\[
2-x=-(x-2),
\]
то
\[
y=-x^2-1,
\qquad
x\ne2.
\]

Выколота точка
\[
A(2;-5).
\]

Ищем значения $k$, при которых прямая
\[
y=kx
\]
имеет с исходным графиком ровно одну общую точку.

\textbf{Случай 1: касание.}

\[
-x^2-1=kx.
\]

Перенесём всё в одну часть:
\[
x^2+kx+1=0.
\]

Для одного корня необходимо:
\[
D=0.
\]

Дискриминант:
\[
D=k^2-4.
\]

Следовательно,
\[
k^2-4=0,
\]
\[
k=2
\qquad
\text{или}
\qquad
k=-2.
\]

\textbf{Случай 2: одна из двух точек пересечения выколота.}

Прямая должна проходить через точку
\[
A(2;-5).
\]

Подставляем её координаты в
\[
y=kx:
\]
\[
-5=2k.
\]

Следовательно,
\[
k=-\frac52.
\]

При этом уравнение пересечения имеет корни
\[
x=2
\qquad\text{и}\qquad
x=\frac12.
\]

Значение
\[
x=2
\]
запрещено ОДЗ, поэтому остаётся ровно одна общая точка.

Итак,
\[
\textbf{
$k\in
\left\{
-\dfrac52,
-2,
2
\right\}$.
}
\]

\begin{center}

\begin{tikzpicture}

\begin{axis}[
    width=0.82\linewidth,
    height=92mm,
    axis lines=middle,
    xmin=-4,
    xmax=4,
    ymin=-10,
    ymax=4,
    xlabel={$x$},
    ylabel={$y$},
    samples=240,
    clip=false,
    unit vector ratio=1 1,
    tick label style={font=\small},
    label style={font=\small},
    legend style={
        font=\small,
        draw=none,
        fill=none,
        at={(0.02,0.02)},
        anchor=south west
    }
]

\addplot[
    very thick,
    color=darkaccent,
    domain=-4:4
]
{-x^2-1};
\addlegendentry{$y=-x^2-1$}

\addplot[
    thick,
    color=accent,
    domain=-4:4
]
{2*x};
\addlegendentry{$y=2x$}

\addplot[
    thick,
    dashed,
    color=accent!80!black,
    domain=-4:4
]
{-2*x};
\addlegendentry{$y=-2x$}

\addplot[
    thick,
    densely dotted,
    color=textgray,
    domain=-4:4
]
{-2.5*x};
\addlegendentry{$y=-\frac52x$}

\addplot[
    only marks,
    mark=o,
    mark size=3.2pt,
    very thick,
    color=darkaccent,
    fill=white
]
coordinates {(2,-5)};

\node[
    anchor=west,
    font=\small,
    text=darkaccent
]
at (axis cs:2.1,-5)
{$A(2;-5)$};

\end{axis}

\end{tikzpicture}

\end{center}

% ============================================================
% КРАТКОЕ ПОВТОРЕНИЕ
% ============================================================

\section*{7. Что особенно важно помнить}
\sectionrule

\begin{itemize}

\item
ОДЗ записывается по \emph{исходной} функции и сохраняется после сокращения.

\item
Скобки $x-a$ и $a-x$ отличаются знаком:
\[
a-x=-(x-a).
\]

\item
После сокращения строится знакомый график, затем на нём удаляются запрещённые точки.

\item
Для пересечения графиков приравниваются формулы функций.

\item
Для касания прямой и параболы в квадратном уравнении пересечения выполняется
\[
D=0.
\]

\item
При наличии выколотой точки условие $D=0$ может дать не все ответы: отдельно проверяется прямая, проходящая через выколотую точку.

\item
В мелких вычислениях особенно контролируйте знаки, квадрат отрицательного числа и перенос слагаемых через знак равенства.

\end{itemize}

% ============================================================
% ТРЕНИРОВОЧНЫЙ БЛОК
% ============================================================

\clearpage

\section*{Тренировочный блок. Путь А}
\sectionrule

Выполните задания последовательно. В графических задачах сначала выпишите ОДЗ, затем упростите функцию, постройте основной график и только после этого отметьте выколотые точки.

\begin{enumerate}

\item
Для функции
\[
y=
\frac{(x+2)(x^2-5x+6)}{x-2}
\]
найдите ОДЗ, упростите выражение и укажите координаты выколотой точки.

\item
Постройте график функции
\[
y=
\frac{(x^2+4)(x-1)}{1-x}.
\]

Укажите:

\begin{itemize}
\item вершину параболы;
\item пересечения с осями;
\item выколотую точку.
\end{itemize}

\item
Для функции из задания 1 определите количество общих точек её графика с прямой
\[
y=m
\]
в зависимости от значения $m$.

\item
Определите все значения $k$, при которых прямая
\[
y=kx
\]
имеет ровно одну общую точку с графиком функции
\[
y=
\frac{(x^2+1)(x-3)}{3-x}.
\]

\item
Определите все значения $k$, при которых прямая
\[
y=kx
\]
имеет ровно одну общую точку с графиком функции
\[
y=
\frac{(x^2+9)(x+2)}{-x-2}.
\]

\end{enumerate}

% ============================================================
% ОТВЕТЫ
% ============================================================

\clearpage

\section*{Ответы}
\sectionrule

\renewcommand{\arraystretch}{1.35}

\begin{table}[h]
\centering

\begin{tabularx}{\linewidth}{@{}>{\bfseries}c X@{}}

\toprule
№ & Ответ \\
\midrule

1 &
$x\ne2$;
\[
y=x^2-x-6;
\]
выколота точка
\[
(2;-4).
\]
\\

2 &
$x\ne1$;
\[
y=-x^2-4.
\]
Вершина:
\[
(0;-4).
\]
С осью $Ox$ пересечений нет.
С осью $Oy$:
\[
(0;-4).
\]
Выколота точка:
\[
(1;-5).
\]
\\

3 &
$0$ точек при
\[
m<-\frac{25}{4};
\]
$1$ точка при
\[
m=-\frac{25}{4};
\]
$1$ точка при
\[
m=-4;
\]
$2$ точки при
\[
m>-\frac{25}{4},
\qquad
m\ne-4.
\]
\\

4 &
\[
k\in
\left\{
-\frac{10}{3},
-2,
2
\right\}.
\]
\\

5 &
\[
k\in
\left\{
-6,
6,
\frac{13}{2}
\right\}.
\]
\\

\bottomrule

\end{tabularx}

\end{table}

\vfill

{\small\color{darkaccent}
Самопроверка: в задачах с параметром убедитесь, что каждый найденный коэффициент действительно оставляет у исходных графиков ровно одну общую точку после учёта ОДЗ.
}

\end{document}